\documentclass[../../main.tex]{subfiles}
\begin{document}
\section{The mass martingale and the stopped centroid reduction}
\label{sec:mass-martingale}

\subsection{Quadratic variation and information rate}

By \eqref{eq:loc-martingale},
\[
  \dd p_t=v_t\cdot\dd W_t,
  \qquad
  v_t=\int_E(x-a_t)\dd\mu_t(x).
\]
Since $a_t=p_tm_t^E+q_tm_t^F$, one has
\begin{equation}\label{eq:v-s-delta}
  v_t=s_t\delta_t .
\end{equation}
Consequently,
\begin{equation}\label{eq:qv-p}
  \dd[p]_t=\abs{v_t}^2\dd t=s_t^2\abs{\delta_t}^2\dd t=s_tr_t\dd t.
\end{equation}
This is the first cut-specific identity: the martingale volatility of the mass is precisely the
centroid gap of the two posterior colors.

Let
\[
  \mathsf h(p)=-p\log p-(1-p)\log(1-p)
\]
be binary entropy. Since $\mathsf h''(p)=-1/(p(1-p))$, Ito's formula and \eqref{eq:qv-p} give
\begin{equation}\label{eq:entropy-rate}
  \frac{\dd}{\dd t}\E\,\mathsf h(p_t)=-\frac12\E r_t .
\end{equation}
Therefore
\begin{equation}\label{eq:total-info}
  \int_0^\infty \E r_t\dd t\le2\mathsf h(p_0)\le2\log2 .
\end{equation}
This exact information identity says that localization gradually reveals the binary label
$\one_E(X)$. It is too weak for KLS: the total information may be bounded while a large initial
spike forces the cut to leave the balanced window rapidly.

\subsection{Intrinsic covariance control}

For any probability measure $\nu$, any set $E$ of mass $p$, and $F=E^c$ of mass $q$, covariance
decomposition gives
\begin{equation}\label{eq:cov-decomp}
  A=p\Sigma_E+q\Sigma_F+pq\delta\delta^T.
\end{equation}
Thus
\begin{equation}\label{eq:B-le-A}
  B:=pq\delta\delta^T\preceq A.
\end{equation}
In particular, on the support of $A$,
\begin{equation}\label{eq:intrinsic-bound}
  pq\,\delta^TA^{-1}\delta\le1.
\end{equation}
The centroid separation is always controlled in the covariance metric of the posterior. KLS needs
Euclidean control; the danger is alignment of $\delta_t$ with large-eigenvalue directions of $A_t$.

A useful warning follows from \eqref{eq:B-le-A}: since $B_t$ has rank one and eigenvalue $r_t$,
\begin{equation}\label{eq:trivial-lambda}
  \abs{v_t}^2=s_tr_t\le s_t\lmax(A_t)\le\tfrac14\lmax(A_t).
\end{equation}
Hence a constant-time estimate of $\E\int_0^T\lmax(A_t)\dd t$ would already imply KLS by
Doob's inequality. Any noncircular proof of the two-color Carleson estimate must exploit the
cut-specific quantities $G_t,B_t,r_t,D_t$, not first assume constant-time spectral control of
$A_t$. Proposition~\ref{prop:ceiling} sharpens this warning into an exact ceiling: the estimate
$\E\int_0^T\lmax(A_t)\dd t\le(1+\kappa)T$ is not merely sufficient for KLS, it is precisely what a
relative-scale excess propagation estimate would require as input, so the relative scale cannot be
a legitimate intermediate target.

\subsection{Stopped centroid estimate}

\begin{assumption}[Stopped centroid estimate]
\label{ass:stopped-centroid}
There exist universal constants $T_0,C$ such that for every isotropic log-concave $\mu$ and every
measurable set $E$ with $p_0\in[2/5,3/5]$,
\begin{equation}\label{eq:stopped-centroid}
  \E\int_0^{T\wedge\tau}\abs{\delta_t}^2\dd t\le CT,
  \qquad 0<T\le T_0 .
\end{equation}
\end{assumption}

\begin{theorem}[Stopped centroid estimate implies KLS]
\label{thm:centroid-implies-kls}
Assumption~\ref{ass:stopped-centroid} implies a dimension-free lower bound on $h_\mu$ for every
isotropic log-concave $\mu$.
\end{theorem}

\begin{proof}
It is enough to prove a universal boundary lower bound for balanced cuts. We give the proof for
$p_0=1/2$; the nested interval $p_0\in[2/5,3/5]$ changes only constants.

If $\tau\le T$, then $\sup_{t\le T\wedge\tau}\abs{p_t-p_0}\ge1/6$. Doob's $L^2$ inequality and
\eqref{eq:qv-p} yield
\[
  \Prob(\tau\le T)
  \le36\E[p]_{T\wedge\tau}
  =36\E\int_0^{T\wedge\tau}s_t^2\abs{\delta_t}^2\dd t
  \le36CT .
\]
Choose a universal $T\le T_0$ so small that $36CT\le1/2$. Then $\Prob(\tau>T)\ge1/2$, and on this
event $\min(p_T,q_T)\ge1/3$.

The posterior $\mu_T$ is $T$-uniformly log-concave. The Bakry--Emery/Lichnerowicz mechanism and the
standard Gaussian isoperimetric comparison for uniformly log-concave measures give
\[
  \mu_T^+(E)\ge c\sqrt T\min(p_T,q_T),
\]
with universal $c>0$; see \cite{BakryGentilLedoux2014} and the convex-measure comparison framework
of \cite{Bobkov1999LogConcave,Milman2009Isoperimetric}. Averaging and using the supermartingale inequality
\eqref{eq:perimeter-supermartingale},
\[
  \mu^+(E)
  \ge \E\mu_T^+(E)
  \ge c\sqrt T\,\E\min(p_T,q_T)
  \ge c'>0 .
\]
Thus every balanced cut has universal boundary measure. By concavity of the isoperimetric profile,
this is equivalent to a dimension-free Cheeger lower bound.
\end{proof}

The same proof shows that it is enough to establish \eqref{eq:stopped-centroid} for a sequence of
balanced near-minimizers of the isoperimetric profile. This observation is what permits the
near-Cheeger geometric route.
\end{document}
